\section{Step 10 Reconstruct the coefficient code}\label{fir-s10}

Shift \texttt{t6} left by eight and combine it into \texttt{t5}. The stored order is the same little-endian order as the input samples. A Q15 scale does not change byte order or make the raw code unsigned. This edit recovers the numerical pattern; the next edit gives that pattern its signed interpretation.

\Needspace{12\baselineskip}
\subsection{Write and inspect}
Replace the temporary branch and nop in \texttt{coefficient\_join}. Keep surrounding labels and region markers.

\Needspace{5\baselineskip}
\MARSCodeLabel{MIPS assembly}
\begin{minted}[samepage=true]{mars-idea}
    sll $t6, $t6, 8
    or $t5, $t5, $t6
\end{minted}
Use recipe \texttt{s10-1} from the companion construction inputs. Pause at \texttt{filter\_ready} before comparing the partial state.

\subsection{Expected observation}
Coefficient bytes 00 c0 reconstruct unsigned pattern 49152 in \texttt{t5}.

\Needspace{11\baselineskip}
\subsection{Check your understanding}
Does the Q15 scale reverse the stored bytes?

\begin{itemize}
  \item \textbf{A} Yes; fractional values use the high byte first.
  \item \textbf{B} Only for negative coefficients.
  \item \textbf{C} No; the contract still stores the low byte first.
\end{itemize}
